\section{Identifiability via the Fisher information matrix}\label{sec:fim}

This section carries out the Fisher information rank analysis that an
earlier draft flagged as outstanding. We review the definition of the
Fisher information matrix and what it tells us about local
identifiability, derive analytically the three rank-deficient parameter
pairs of the pre-G1 FSA-v2 model, present the reparametrization that
absorbs them, and report the post-reparametrization evidence that all
$30$ parameters of the joint state-space model are now identifiable
from the four-channel observation record at the one-day window scale.

\subsection{Fisher information and local identifiability}

For a parameter vector $\theta \in \R^d$ and an observation record
$y$ from the model, the score function is
$s(\theta;y) = \nabla_\theta \log p(y\mid\theta)$ and the
\emph{Fisher information matrix} is
\begin{equation}
  \Ical(\theta) \;=\;
    \E_{y \sim p(\cdot\mid\theta)}\!\bigl[\,s(\theta;y)\,s(\theta;y)^{\!\top}\bigr].
  \label{eq:fim-def}
\end{equation}
$\Ical(\theta)$ is symmetric positive semidefinite. Two facts make it
the right object to look at when asking whether the data can constrain
$\theta$.

\begin{proposition}[Local identifiability]\label{prop:fim-id}
Suppose $\theta \mapsto p(y\mid\theta)$ is twice differentiable in a
neighbourhood of $\theta^*$. Then $\theta^*$ is locally identifiable if
and only if $\Ical(\theta^*)$ has full rank $d$.
\end{proposition}

\begin{proposition}[Cram\'er--Rao]\label{prop:fim-cr}
Any unbiased estimator $\hat\theta$ of $\theta$ has covariance bounded
below in the positive-semidefinite ordering by
$\Ical(\theta)^{-1}$ (when $\Ical$ is invertible).
\end{proposition}

Proposition~\ref{prop:fim-id} says the FIM rank is the diagnostic for
identifiability; Proposition~\ref{prop:fim-cr} says its inverse sets
the floor for posterior-variance scaling, so a near-singular FIM
implies posteriors with very wide spread along the singular directions.
In this code base we are interested in identifiability at one-day
windows -- the rolling-window unit of the SMC${}^2$ filter
(Section~\ref{sec:smc2}) -- and at a representative truth point near
the operating point of Definition~\ref{def:operating-point}.

\subsection{Analytical FIM rank deficiencies of the pre-G1 model}

The four-channel observation model
\eqref{eq:obs-hr}--\eqref{eq:obs-steps} is, conditional on the latent
state, a product of one-dimensional exponential-family likelihoods
whose score functions are linear in $(B,F,A)$. The state itself is
driven by the drift \eqref{eq:fsa-B}--\eqref{eq:fsa-A}. At a one-day
window the slow Banister state moves only marginally, so the
information about the dynamics parameters is carried by how
$\partial(B,F,A)/\partial\theta$ couples through the observation model.

A direct calculation along these lines reveals three pairs of
parameters whose effects on the one-day data are exactly collinear, and
hence whose contribution to the FIM is rank-one rather than rank-two
within each pair. Recall the pre-G1 drift was
\begin{equation*}
  \dd B = \bigl[\,\kappa_B(1+\varepsilon_A A)\,\Phi - B/\tau_B\,\bigr]\dd t + \cdots,
  \quad
  \dd F = \bigl[\,\kappa_F\,\Phi - (1+\lambda_A A)/\tau_F \cdot F\,\bigr]\dd t + \cdots,
\end{equation*}
and $\mu(B,F) = \mu_0 + \mu_B B - \mu_F F - \mu_{FF} F^2$. At constant
control $\Phi \equiv \Phi_{\mathrm{typ}}$ and near
$A \equiv A_{\mathrm{typ}}$, $F \equiv F_{\mathrm{typ}}$:

\paragraph{Pair 1: $(\kappa_B,\varepsilon_A)$.} Linearising
$1+\varepsilon_A A$ around $A_{\mathrm{typ}}$ and treating fluctuations
$\delta A = A - A_{\mathrm{typ}}$ as small at one-day windows, the
$B$-drift is
\begin{equation}
  \kappa_B (1+\varepsilon_A A)\,\Phi
   \;\approx\;
   \underbrace{\kappa_B(1+\varepsilon_A A_{\mathrm{typ}})}_{=:\,\kappa_B^{\mathrm{eff}}}\,\Phi
   \;+\; \kappa_B\,\varepsilon_A\,\delta A\,\Phi.
  \label{eq:pair1-expand}
\end{equation}
The first term carries data-relevant information through the product
$\kappa_B(1+\varepsilon_A A_{\mathrm{typ}})$, and is identifiable. The
second term is $O(\delta A)$, vanishing at the operating point and
small in one-day windows; the data therefore see $\kappa_B$ and
$\varepsilon_A$ only through the combination $\kappa_B^{\mathrm{eff}}$,
making the pair rank-deficient.

\paragraph{Pair 2: $(\mu_F,\mu_{FF})$.} Expanding the quadratic
$\mu_{FF} F^2$ around $F_{\mathrm{typ}}$,
\begin{align*}
  -\mu_F F - \mu_{FF} F^2
  &= -\mu_F F - \mu_{FF}\bigl[(F-F_{\mathrm{typ}})^2 + 2 F_{\mathrm{typ}} F - F_{\mathrm{typ}}^2\bigr]\\
  &= -\bigl(\,\underbrace{\mu_F + 2 F_{\mathrm{typ}} \mu_{FF}}_{=:\,\mu_F^{\mathrm{eff}}}\,\bigr) F
     - \mu_{FF}(F-F_{\mathrm{typ}})^2 + \mu_{FF} F_{\mathrm{typ}}^2.
\end{align*}
At one-day windows $\delta F = F - F_{\mathrm{typ}}$ is small, so the
$(F-F_{\mathrm{typ}})^2$ term is second-order and uninformative. The
data therefore see $\mu_F$ and $\mu_{FF}$ only through the linear slope
$\mu_F^{\mathrm{eff}}$, giving a second rank-deficient pair.

\paragraph{Pair 3: $(\tau_F,\lambda_A)$.} The same linearisation
applied to the $F$-clearance factor gives
\begin{equation*}
  \frac{1+\lambda_A A}{\tau_F}
  \;\approx\;
  \frac{1+\lambda_A A_{\mathrm{typ}}}{\tau_F}
  + \frac{\lambda_A}{\tau_F}\,\delta A
  \;=:\;
  \frac{1}{\tau_F^{\mathrm{eff}}} + O(\delta A).
\end{equation*}
The data therefore see $\tau_F$ and $\lambda_A$ only through the
effective rate $1/\tau_F^{\mathrm{eff}} = (1+\lambda_A A_{\mathrm{typ}})/\tau_F$,
giving the third rank-deficient pair.

\smallskip
At a one-day truth window centred on the operating point, the
Fisher information matrix of the pre-G1 model therefore has rank
$d - 3$. The three null directions correspond exactly to the
combinations
\begin{equation}
  (\delta\kappa_B,\delta\varepsilon_A) \perp \nabla \kappa_B^{\mathrm{eff}},
  \quad
  (\delta\mu_F,\delta\mu_{FF}) \perp \nabla \mu_F^{\mathrm{eff}},
  \quad
  (\delta\tau_F,\delta\lambda_A) \perp \nabla(1/\tau_F^{\mathrm{eff}}),
  \label{eq:fim-nulls}
\end{equation}
i.e.\ to perturbations along which the strongly identified linear
combination is held fixed.

\subsection{The G1 reparametrization absorbs the three pairs}

The cure is straightforward: rotate each rank-deficient pair into a
\emph{strongly identified} linear combination and a
\emph{weakly identified residual} that vanishes at the operating
point. The G1 reparametrization, which we already wrote out as
equations \eqref{eq:fsa-mu}--\eqref{eq:fsa-A}, does exactly this. A
side-by-side comparison is in Table~\ref{tab:g1-reparam}.

\begin{table}[h]
\centering
\small
\begin{tabular}{@{}llll@{}}
\toprule
Pair & Strongly identified & Residual & Comment \\
\midrule
$(\kappa_B,\varepsilon_A)$  & $\kappa_B^{\mathrm{new}} = \kappa_B(1+\varepsilon_A A_{\mathrm{typ}})$
                             & $\varepsilon_A$  & residual factor $(1+\varepsilon_A A)/(1+\varepsilon_A A_{\mathrm{typ}}) \to 1$ at $A_{\mathrm{typ}}$ \\
$(\mu_F,\mu_{FF})$          & $\mu_F^{\mathrm{new}} = \mu_F + 2 F_{\mathrm{typ}} \mu_{FF}$
                             & $\mu_{FF}$       & curvature on $(F - F_{\mathrm{typ}})^2 \to 0$ at $F_{\mathrm{typ}}$ \\
$(\tau_F,\lambda_A)$        & $\tau_F^{\mathrm{new}} = \tau_F/(1+\lambda_A A_{\mathrm{typ}})$
                             & $\lambda_A$      & residual factor $(1+\lambda_A A)/(1+\lambda_A A_{\mathrm{typ}}) \to 1$ at $A_{\mathrm{typ}}$ \\
\bottomrule
\end{tabular}
\caption{The Stage~G1 reparametrization. Each rank-deficient parameter
pair of the pre-G1 model is rotated into a strongly identified
combination (left column) and a residual (middle column) that vanishes
at the operating point of Definition~\ref{def:operating-point}. The
parameter \emph{names} are kept; only the meanings and truth values
change. There is also an implicit shift in $\mu_0$ that absorbs the
constant from completing the square, $\mu_0 \mapsto \mu_0 + \mu_{FF}
F_{\mathrm{typ}}^2$.}
\label{tab:g1-reparam}
\end{table}

\begin{proposition}[G1 is mathematically equivalent to v2]\label{prop:g1-equiv}
The G1 drift \eqref{eq:fsa-B}--\eqref{eq:fsa-A} with truth values from
Table~\ref{tab:fsa-truth} is identically equal to the pre-G1 drift
under the relabelling
\begin{align*}
  \kappa_B^{\mathrm{old}} &= \kappa_B/(1+\varepsilon_A A_{\mathrm{typ}}), &
  \tau_F^{\mathrm{old}}   &= \tau_F\,(1+\lambda_A A_{\mathrm{typ}}), \\
  \mu_F^{\mathrm{old}}    &= \mu_F - 2 F_{\mathrm{typ}} \mu_{FF}, &
  \mu_0^{\mathrm{old}}    &= \mu_0 - \mu_{FF} F_{\mathrm{typ}}^2.
\end{align*}
\end{proposition}

The equivalence is verified by direct substitution; the project's unit
test \texttt{tests/test\_g1\_reparam.py} performs the same check
numerically over $50$ random states and confirms the two drifts agree
to floating-point precision.

\subsection{Empirical evidence of identifiability after G1}\label{sec:fim-empirical}

The analytical argument above shows that the G1 form has full FIM rank
at the operating point. The implementation furnishes empirical
evidence of the same fact at one-day rolling windows on simulated
truth data.

\begin{table}[h]
\centering
\small
\begin{tabular}{@{}lcc@{}}
\toprule
Quantity & Pre-G1 & Post-G1 \\
\midrule
$\rank \widehat{\Ical}$ at one-day window               & $d-3$            & $d$ \\
parameters covered by 90\,\% credible interval   & $27/30$          & $\mathbf{30/30}$ \\
identifiable estimable parameters                & $3/6$            & $\mathbf{6/6}$ \\
posterior std on $\varepsilon_A$ (lognormal scale) & $0.093$          & $\mathbf{0.016}$ \\
$90\,\%$ CI for $\varepsilon_A$                  & $[0.275,\,0.550]$ & $\mathbf{[0.373,\,0.423]}$ \\
\bottomrule
\end{tabular}
\caption{Identifiability before and after the G1 reparametrization,
measured at one-day rolling-window scale on simulated truth data.
``Estimable'' counts the six parameters that the FIM analysis
identified as candidates for tightening. Numbers from
\texttt{version\_2/outputs/fsa\_high\_res/} and the Stage~G1 commit
record. The residual standard deviation on $\varepsilon_A$ tightens by
roughly a factor of six; analogous tightening is observed on the other
two residuals $\mu_{FF}$ and $\lambda_A$.}
\label{tab:g1-evidence}
\end{table}

The empirical findings reported in Table~\ref{tab:g1-evidence} support
the analytical claim of Section~\ref{sec:fim} above: the G1 form is
identifiable at one-day windows, the strongly identified parameters
$(\kappa_B,\tau_F,\mu_F)$ retain their original prior widths, and the
residuals $(\varepsilon_A,\lambda_A,\mu_{FF})$ -- whose posteriors were
formerly artificially wide because of the rank deficiency -- now
contract to roughly a sixth of their pre-G1 spread under tightened
priors that the rank analysis itself motivates.

\begin{remark}[Where the FIM analysis lives]\label{rem:fim-where}
The analytical FIM derivation summarised here is recorded in the
project's commit log (Stage~G1) rather than as standalone code. A
useful follow-up would be to add a numerical FIM module under
\texttt{smc2fc/diagnostics/fim/} that reproduces these results from the
JAX dynamics by automatic differentiation, so that any future
modification of the model triggers a re-run of the rank check
automatically.
\end{remark}
